% ==================== 1.6. TENDENCIAS Y LÍNEAS FUTURAS ====================
        \section{Tendencias y líneas futuras}
            El análisis de los estudios revisados pone sobre la mesa que la Inteligencia Artificial aplicada al análisis de la voz para la detección temprana de enfermedades neurodegenerativas es un campo en expansión, con resultados prometedores pero que aún está en una fase experimental. A partir de las evidencias recopiladas, pueden observarse varias tendencias tecnológicas y líneas de investigación futura que marcan el paso de la evolución de esta disciplina.

            \subsection{Avances en modelos de representación y foundation models}
                Una de las líneas más claras es la transición desde modelos basados en features manuales hacia representaciones profundas aprendidas automáticamente. Los estudios más recientes (Agbavor \& Liang, 2022; Dao et al., 2025) demuestran el potencial de los modelos auto-supervisados como: data2vec, wav2vec2.0 o Whisper, capaces de aprender representaciones acústicas sin necesidad de anotaciones.
                Estos foundation models permiten procesar grandes volúmenes de datos no etiquetados y facilitan la transferencia entre tareas, reduciendo la dependencia de datasets específicos. En el futuro, el uso de estos modelos preentrenados mejorará la generalización, acelerará el desarrollo de nuevos sistemas y sentará las bases para herramientas diagnósticas basadas en voz.

            \subsubsection{Integración multimodal de datos}
                Varios estudios coinciden en la necesidad de combinar información acústica, prosódica, lingüística y textual para mejorar la precisión diagnóstica, especialmente en enfermedades donde los síntomas cognitivos son más marcados, como el Alzheimer (García-Gutiérrez et al., 2024; He et al., 2023).
                El análisis conjunto de voz y texto, o de voz y datos clínicos complementarios, constituye una tendencia ascendente hacia modelos multimodales capaces de reflejar de forma más completa los cambios neurológicos. La fusión de estas fuentes de información permitirá desarrollar sistemas más robustos y sensibles a las distintas manifestaciones de la enfermedad.

            \subsubsection{Explicabilidad y validación clínica}
                Los avances técnicos deben ir acompañados de una mayor transparencia e interpretabilidad de los modelos. Las aproximaciones recientes que incorporan Explainable AI (Shen et al., 2024; Muntaha et al., 2024) enseñan que es posible identificar qué características vocales influyentes en la clasificación, lo que facilita la validación por parte de especialistas.
                En un futuro próximo se espera un incremento de los trabajos centrados en IA explicable y validada clínicamente, condición necesaria para la incorporación de estos sistemas en entornos sanitarios. La colaboración entre ingenieros, lingüistas y médicos será la clave para traducir los resultados de laboratorio en herramientas diagnósticas reales.

            \subsubsection{Creación de bases de datos abiertas y estandarizadas}
                La falta de datasets grandes y uniformes es una de las principales barreras actuales. La tendencia internacional apunta hacia la construcción de bases de datos abiertas, multilingües y balanceadas, que permitan comparar modelos bajo las mismas condiciones experimentales. Iniciativas de este tipo facilitarían la reproducibilidad de los resultados y el desarrollo de modelos robustos y equitativos.
                Además, la estandarización de protocolos de grabación y tareas vocales (vocales sostenidas, lectura, discurso libre) será esencial para reducir la variabilidad y establecer puntos de referencia comunes en la comunidad científica.
       
            \subsubsection{Aplicaciones en entornos reales y monitorización continua}
                Los avances en hardware y en conectividad han abierto la posibilidad de aplicar estos sistemas fuera del laboratorio. Estudios como el de Worasawate et al. (2021) muestran la viabilidad de utilizar grabaciones obtenidas con smartphones para detectar alteraciones vocales, por lo que se prevee un futuro donde las herramientas de cribado puedan implementarse en dispositivos de uso cotidiano.
                La tendencia apunta hacia la monitorización continua y remota, permitiendo un seguimiento no invasivo y económico, así como la detección temprana de recaídas o empeoramientos.
       
            \subsubsection{Síntesis prospectiva}
                En conjunto, las líneas futuras del campo pueden resumirse en cinco ejes principales:

                \begin{itemize}
                    \item \textbf{Automatización y auto-supervisión}, adoptando foundation models para extraer representaciones generales del habla.
                    \item \textbf{Multimodalidad}, realizando la fusión de la voz, texto y otros datos clínicos.
                    \item \textbf{Explicabilidad}, llevando a cabo una incorporación sistemática de técnicas XAI para aumentar la confianza clínica.
                    \item \textbf{Estandarización de datos}, creando datasets abiertos, equilibrados y comparables.
                    \item \textbf{Aplicación práctica}, implementando plataformas móviles y entornos reales para monitorización temprana.
                \end{itemize}

                Estas tendencias apuntan hacia un mundo donde la voz se asienta como un biomarcador accesible y fiable, y donde la Inteligencia Artificial, combinada con una validación clínica, permite el desarrollo de sistemas de detección precoz con impacto real en la práctica médica.
       
       